\documentclass[10pt,a4paper]{article}
\usepackage[a4paper,margin=14mm,top=12mm,bottom=14mm]{geometry}
\usepackage{fontspec}
\setmainfont{Liberation Sans}
\usepackage{polyglossia}
\setdefaultlanguage{french}
\usepackage{amsmath}
\usepackage{xcolor}
\usepackage{graphicx}
\usepackage[most]{tcolorbox}
\usepackage{tabularx,booktabs,array,multirow}
\usepackage{enumitem}
\usepackage{fancyhdr}
\usepackage{lastpage}
\usepackage[hidelinks]{hyperref}
\definecolor{uimm}{HTML}{005677}
\definecolor{cyan}{HTML}{13A2DD}
\definecolor{fond}{HTML}{E6F2FA}
\definecolor{ok}{HTML}{2E7D32}
\definecolor{ko}{HTML}{B3261E}
\definecolor{sylvain}{HTML}{7B1FA2}
\setlength{\parindent}{0pt}
\setlength{\parskip}{3pt}
\setlist{nosep}
\setlist[itemize]{leftmargin=13pt,label=\textcolor{cyan}{\textbullet}}
\newcommand{\titre}[1]{\vspace{5pt}{\color{uimm}\large\bfseries #1}\par\vspace{-3pt}{\color{cyan}\rule{\linewidth}{1.2pt}}\par\vspace{0pt}}
\tcbset{boite/.style={colback=fond,colframe=uimm,boxrule=0.6pt,arc=2pt,left=6pt,right=6pt,top=4pt,bottom=4pt},
        defi/.style={colback=sylvain!6,colframe=sylvain,boxrule=0.8pt,arc=2pt,left=6pt,right=6pt,top=4pt,bottom=4pt,
                     title={\bfseries Défi},coltitle=white,colbacktitle=sylvain,fonttitle=\small}}
\newcommand{\bandeau}[3]{%
\begin{tcolorbox}[colback=uimm,colframe=uimm,arc=2pt,left=4pt,right=8pt,top=4pt,bottom=4pt]
\begin{minipage}[c]{18mm}\setlength{\fboxsep}{1.2pt}\fcolorbox{white}{white}{\includegraphics[height=16.5mm]{logo_uimm.jpg}}\end{minipage}\hspace{4mm}%
\begin{minipage}[c]{\dimexpr\linewidth-22mm\relax}\color{white}
{\small #1}\\[1pt]{\LARGE\bfseries #2}\\[2pt]{\large #3}\hfill{\small CQPM Technicien qualité}
\end{minipage}
\end{tcolorbox}}
\newcommand{\reponse}[1][1]{\par\vspace{2pt}{\color{gray!60}\foreach \i in {1,...,#1}{\rule{\linewidth}{0.3pt}\par\vspace{9pt}}}}
\usepackage{pgffor}
\pagestyle{fancy}\fancyhf{}
\renewcommand{\headrulewidth}{0pt}
\newcommand{\pied}[1]{\fancyfoot[L]{\footnotesize\color{uimm}#1}\fancyfoot[R]{\footnotesize\color{uimm}Page \thepage{} sur \pageref{LastPage}}}

\usepackage{tikz}
\usepackage{colortbl}
\pied{TP R\&R · L'opérateur compte-t-il ? · fiche stagiaire}
\begin{document}
\thispagestyle{fancy}
\bandeau{Travaux pratiques · Capabilité du moyen de mesure}{L'opérateur compte-t-il ?}{Étude R\&R (répétabilité et reproductibilité), méthode Bosch}

\begin{tcolorbox}[boite]
\textbf{La question.} Trois personnes mesurent les mêmes butées avec le même micromètre. Trouvent-elles la même chose ?
Si non, l'écart vient-il de \textbf{l'instrument} (une même personne ne retrouve pas sa valeur) ou \textbf{des opérateurs} (chacun mesure à sa façon) ?
\end{tcolorbox}

\titre{1. Le GRR en quatre images}
\begin{minipage}[t]{0.70\linewidth}
\begin{enumerate}[label=\textcolor{white}{\colorbox{uimm}{\textbf{\arabic*}}},leftmargin=20pt,itemsep=3pt]
\item \textbf{Est-ce que je me répète ?} Je mesure deux fois la même butée : l'écart entre mes deux valeurs, c'est la \textbf{répétabilité} (EV). Elle dépend surtout de l'instrument et de la façon de le poser.
\item \textbf{Sommes-nous d'accord ?} On compare la moyenne de chaque opérateur : l'écart entre ces moyennes, c'est la \textbf{reproductibilité} (AV). Elle vient des méthodes de chacun.
\item \textbf{Comment additionner ?} Les deux causes s'assemblent comme les côtés d'un \textbf{triangle rectangle} : $GRR^2 = EV^2 + AV^2$. Les pourcentages ne s'ajoutent donc pas.
\item \textbf{Qu'est-ce que ça change au poste ?} Une valeur lue n'est pas un point, c'est un intervalle d'environ $\pm\,3\,GRR$. Près des limites de tolérance, l'instrument ne sait plus trancher : c'est la \textbf{zone de doute}.
\end{enumerate}
\end{minipage}\hfill
\begin{minipage}[t]{0.27\linewidth}\vspace{4pt}\centering
\begin{tikzpicture}[scale=0.95]
\fill[fond] (0,0)--(3,0)--(3,1.8)--cycle;
\draw[line width=2.2pt,color=cyan!70!blue] (0,0)--(3,0) node[midway,below,font=\scriptsize]{EV (instrument)};
\draw[line width=2.2pt,color=orange!80!black] (3,0)--(3,1.8) node[midway,right,font=\scriptsize,align=left]{AV\\(opérateurs)};
\draw[line width=2.2pt,color=uimm] (0,0)--(3,1.8) node[midway,above left,font=\scriptsize\bfseries]{GRR};
\draw (2.8,0)--(2.8,0.2)--(3,0.2);
\end{tikzpicture}\\[6pt]
{\small Le \%GRR compare le flou de la mesure à la tolérance :}
\[\%GRR=\frac{6\,GRR}{T}\times100\ \%\]
\end{minipage}

\renewcommand{\arraystretch}{1.2}
\begin{tabularx}{\linewidth}{@{}l X@{}}
\toprule
\textcolor{ok}{\textbf{\%GRR $\le$ 10 \%}} & Système de mesure \textbf{capable}.\\
\textbf{10 \% $<$ \%GRR $\le$ 30 \%} & \textbf{Capable sous conditions} : utilisable en acceptant un risque, à améliorer.\\
\textcolor{ko}{\textbf{\%GRR $>$ 30 \%}} & \textbf{Non capable}.\\
\bottomrule
\end{tabularx}

\titre{2. Le matériel et les rôles}
\begin{itemize}
\item 10 butées numérotées de 1 à 10, mesurées au \textbf{trait repère}. Un micromètre (résolution 0,001 mm). Tolérance de la butée : 11,85 à 12,00 mm, soit $T = 0{,}15$ mm.
\item \textbf{3 opérateurs} (A, B, C) mesurent. \textbf{2 secrétaires} : le premier donne les butées dans l'ordre tiré au sort, le second note les valeurs.
\item La page web \textbf{« L'opérateur compte-t-il ? »} : elle tire l'ordre au sort et fait tous les calculs.
\end{itemize}

\titre{3. Le protocole}
\begin{enumerate}[label=\textcolor{cyan}{\textbf{\arabic*.}},leftmargin=15pt]
\item Sur la page, cliquer sur « Tirer l'ordre de mesure au sort ». Recopier les six ordres page 2.
\item \textbf{Série 1.} L'opérateur A mesure les 10 butées dans son ordre. Puis B, puis C. L'opérateur \textbf{ne voit pas} ce que le secrétaire note.
\item \textbf{Série 2.} On recommence seulement quand les trois ont fini la série 1, chacun avec son nouvel ordre.
\item Toujours au trait repère, même serrage (cliquet du micromètre). Noter toutes les valeurs, \textbf{n'en supprimer aucune}.
\item \textbf{Avant de calculer, à la main} (grille page 2) : chaque opérateur calcule ses 10 écarts entre série 1 et série 2 et \textbf{entoure le plus grand} (image 1). Puis on calcule la moyenne de chaque opérateur et on les compare (image 2).
\item Saisir les 60 valeurs dans la page, cliquer sur « Calculer », puis ouvrir le volet « Comprendre d'où vient le GRR ».
\end{enumerate}
{\small\textbf{Pourquoi tirer au sort ?} Si l'opérateur reconnaît la pièce, il risque de retrouver sa première valeur de mémoire : on croirait l'instrument meilleur qu'il n'est.}

\newpage
\titre{4. L'ordre de passage (à recopier depuis la page)}
\renewcommand{\arraystretch}{1.45}
\begin{tabularx}{\linewidth}{|l|*{10}{>{\centering\arraybackslash}X|}}
\hline
\textbf{Opérateur · série} & \textbf{1\textsuperscript{re}} & \textbf{2\textsuperscript{e}} & \textbf{3\textsuperscript{e}} & \textbf{4\textsuperscript{e}} & \textbf{5\textsuperscript{e}} & \textbf{6\textsuperscript{e}} & \textbf{7\textsuperscript{e}} & \textbf{8\textsuperscript{e}} & \textbf{9\textsuperscript{e}} & \textbf{10\textsuperscript{e}}\\\hline
A · série 1 & & & & & & & & & & \\\hline
B · série 1 & & & & & & & & & & \\\hline
C · série 1 & & & & & & & & & & \\\hline
A · série 2 & & & & & & & & & & \\\hline
B · série 2 & & & & & & & & & & \\\hline
C · série 2 & & & & & & & & & & \\\hline
\end{tabularx}

\titre{5. La grille de relevé (mm)}
\renewcommand{\arraystretch}{1.6}
\begin{tabularx}{\linewidth}{|c|*{9}{>{\centering\arraybackslash}X|}}
\hline
\multirow{2}{*}{\textbf{Butée}} & \multicolumn{3}{c|}{\textbf{Opérateur A :} \rule{16mm}{0.3pt}} & \multicolumn{3}{c|}{\textbf{Opérateur B :} \rule{16mm}{0.3pt}} & \multicolumn{3}{c|}{\textbf{Opérateur C :} \rule{16mm}{0.3pt}}\\
\cline{2-10}
 & Série 1 & Série 2 & \textbf{Écart} & Série 1 & Série 2 & \textbf{Écart} & Série 1 & Série 2 & \textbf{Écart}\\\hline
1 & & & & & & & & & \\\hline
2 & & & & & & & & & \\\hline
3 & & & & & & & & & \\\hline
4 & & & & & & & & & \\\hline
5 & & & & & & & & & \\\hline
6 & & & & & & & & & \\\hline
7 & & & & & & & & & \\\hline
8 & & & & & & & & & \\\hline
9 & & & & & & & & & \\\hline
10 & & & & & & & & & \\\hline
\textbf{Moyenne} & \multicolumn{2}{c|}{} & & \multicolumn{2}{c|}{} & & \multicolumn{2}{c|}{} & \\\hline
\end{tabularx}

\vspace{3pt}{\small Ligne « Moyenne » : sous les séries, la moyenne des 20 mesures de l'opérateur (image 2) ; dans la colonne « Écart », la moyenne de ses 10 écarts (image 1).}

\begin{tcolorbox}[boite]
\textbf{À savoir.} Le livret Bosch impose une étude de type 1 (Cg, Cgk sur un étalon) \textbf{avant} le R\&R.
Un R\&R non capable n'accuse pas forcément l'instrument : des pièces irrégulières (diamètre qui varie selon l'endroit mesuré) donnent le même symptôme. D'où le trait repère.
\end{tcolorbox}

\newpage
\titre{6. Les questions}
\begin{enumerate}[label=\textcolor{cyan}{\textbf{Q\arabic*.}},leftmargin=22pt]
\item La résolution du micromètre représente quel pourcentage de $T$ ? Est-ce acceptable (moins de 5 \%) ?\reponse
\item Quel est le \%GRR ? Le système de mesure est-il capable, capable sous conditions, ou non capable ?\reponse
\item Quelle part est la plus grande : celle de l'instrument (\%EV) ou celle des opérateurs (\%AV) ? Votre calcul à la main (écarts et moyennes, page 2) le laissait-il prévoir ?\reponse[2]
\item Sur le graphique « Moyenne de chaque opérateur », un opérateur se détache-t-il ? Quelle explication voyez-vous ?\reponse[2]
\item Le \%EV et le \%AV ajoutés donnent plus que le \%GRR. Pourquoi ? (image 3)\reponse[2]
\item Une butée est lue à 11,996 mm. Avec votre GRR, peut-on être sûr qu'elle est dans la tolérance ? (image 4)\reponse[2]
\item Proposez une action concrète pour faire baisser le \%GRR.\reponse[2]
\end{enumerate}

\begin{tcolorbox}[defi]
\textbf{Calcul à la main par la méthode des étendues moyennes} (livret, annexe D.3). Le livret la juge dépassée et recommande l'analyse de variance faite par la page, mais elle montre l'idée.
\begin{enumerate}[label=\textbf{\alph*.},leftmargin=14pt]
\item Faire la moyenne des 30 écarts de la grille : $\bar R$.
\item Répétabilité : $EV = 0{,}8862 \times \bar R$.
\item Prendre les trois moyennes d'opérateurs et l'écart $R_{\bar x}$ entre la plus grande et la plus petite.
\item Reproductibilité : $AV=\sqrt{(0{,}5231\times R_{\bar x})^2 - EV^2/20}$ (si le nombre sous la racine est négatif, $AV = 0$).
\item $GRR=\sqrt{EV^2+AV^2}$, puis \%GRR. Comparer avec la page. Pourquoi les deux résultats diffèrent-ils ?
\end{enumerate}
\end{tcolorbox}

\vfill
{\footnotesize\color{uimm}\textbf{Source.} Robert Bosch GmbH, \emph{Booklet 10 : Capability of Measurement and Test Processes}, édition novembre 2019 : § 4.2 (procédure 2 : au moins 10 pièces, 3 opérateurs, 2 séries ; seuils du \%GRR), annexes D.2 (analyse de variance), D.3 et D.4 (étendues moyennes, facteurs $K_1 = 0{,}8862$ et $K_2 = 0{,}5231$).}
\end{document}
